% Zdroj: PDF priklady-jaro-2024.pdf, fyzická str. 22. Značka: specializace UI-SU. Zařazeno: S3.
\section{Shlukování}
\szzotazka{jaro 2024}{28}{Shlukování}

\noindent\textbf{Zadání.}\par
Máme data zadaná v tabulce níže a graficky znázorněná v obrázku. Chtěli bychom je rozdělit do třech shluků pomocí
algoritmu $k$-means.

\medskip
\noindent\textit{Poznámka:} Jména nejsou součástí dat -- slouží jen pro vizualizaci a pro zjednodušení formulace odpovědí.

\begin{center}
\begin{minipage}{0.32\linewidth}
\centering
\begin{tabular}{c|cc}
\toprule
Jméno & x & y \\
\midrule
A & $-2$ & $1$ \\
B & $-1$ & $-1$ \\
C & $3$ & $-2$ \\
D & $4$ & $-1$ \\
E & $3$ & $3$ \\
F & $4$ & $4$ \\
G & $-3$ & $0$ \\
H & $4$ & $2$ \\
I & $4$ & $-2$ \\
J & $6$ & $1$ \\
\bottomrule
\end{tabular}
\end{minipage}\hfill
\begin{minipage}{0.62\linewidth}
\centering
\includegraphics[width=\linewidth]{shlukovani-jaro2024.png}
\end{minipage}
\end{center}

\begin{enumerate}
  \item Popište stručně hlavní kroky algoritmu $k$-means.
  \item Předpokládejme, že v první iteraci jsme náhodně vybrali body $A$, $F$ a $J$ jako počáteční středy shluků. Jaké bude
    rozdělení bodů do shluků po dokončení první iterace? Jaké budou středy těchto shluků? Při výpočtech používejte
    Euklidovskou vzdálenost.
  \item Změní se nějak středy shluků v dalších iteracích algoritmu?
\end{enumerate}

\medskip
\noindent\textbf{Náčrt řešení.}\par
\begin{enumerate}
  \item Algoritmus napřed náhodně zvolí počáteční středy shluků. Následně opakuje dva kroky: přiřazení každého bodu k
    nejbližšímu středu shluku a přepočítání polohy středu (jako průměr bodů patřících do daného shluku).
  \item V prvním shluku budou body $\{A, G, B\}$, ve druhém shluku budou body $\{E, F, H\}$, ve třetím budou body $\{J, D, I, C\}$.
    Nové středy budou průměry těchto bodů $(-2, 0)$, $(\tfrac{11}{3}, 3)$ a $(\tfrac{17}{4}, -1)$.
  \item Přiřazení do shluků už se v další iteraci nezmění, středy shluků tedy také ne.
\end{enumerate}

\medskip
\noindent\textit{Doplňující vysvětlení (nad rámec oficiálního náčrtu):} Hraniční body: $C(3,-2)$ má ke středům $A,F,J$ vzdálenosti $\sqrt{34},\sqrt{37},\sqrt{18}$, patří tedy k~$J$; $D(4,-1)$ má vzdálenosti $\sqrt{40},5,\sqrt8$, také $J$; $H(4,2)$ má k~$F$ vzdálenost $2$ a~k~$J$ jen $\sqrt5$, patří k~$F$. Ve druhé iteraci ke středům $(-2,0)$, $(\tfrac{11}{3},3)$, $(\tfrac{17}{4},-1)$ vyjde stejné přiřazení (ověřeno výpočtem), takže algoritmus končí.
